\documentclass[a4paper, 12pt]{article}
\usepackage[portuges]{babel}
\usepackage[utf8]{inputenc}
\usepackage{amsmath}
\usepackage{indentfirst}
\usepackage{graphicx}
\usepackage{multicol,lipsum}

\usepackage{xcolor}
\usepackage{hyperref}
\hypersetup{
    colorlinks=true,
    % linkcolor=cyan,
    filecolor=magenta,      
    urlcolor=blue,
    citecolor=green,
}

\begin{document}
%\maketitle

\begin{titlepage}
	\begin{center}
	
	%\begin{figure}[!ht]
	%\centering
	%\includegraphics[width=2cm]{c:/ufba.jpg}
	%\end{figure}

		\Huge{UDESC}\\
		\large{Departamento de Ciência da Computação}\\ 
		\vspace{15pt}
        \vspace{95pt}
        \textbf{\LARGE{Teoria dos Grafos: Tarefa 1}}\\
		%\title{{\large{Título}}}
		\vspace{3,5cm}
	\end{center}
	
	\begin{flushleft}
		\begin{tabbing}
			Representação Computacional de um grafo através de uma Lista de Adjacências \\ \\
			Alunos: Heitor Linhares Caetano e Deivid Veras Lourenço \\
			Professor: Gilmário Barbosa dos Santos \\
	      \setlength{\labelsep}{\leftmargin}
	\end{tabbing}
 \end{flushleft}
	\vspace{1cm}
	
	\begin{center}
		\vspace{\fill}
			 Maio\\
			 2025
			\end{center}
\end{titlepage}
%%%%%%%%%%%%%%%%%%%%%%%%%%%%%%%%%%%%%%%%%%%%%%%%%%%%%%%%%%%

% % % % % % % % %FOLHA DE ROSTO % % % % % % % % % %

\newpage
% % % % % % % % % % % % % % % % % % % % % % % % % %
\newpage
\tableofcontents
\thispagestyle{empty}

\newpage
\pagenumbering{arabic}
% % % % % % % % % % % % % % % % % % % % % % % % % % %

\section{Objetivo do Trabalho}

A atividade consistiu em implementar um grafo através da representação computacional
\textit{Lista de Adjacências} e que é carregado a partir de uma base de dados no
formato CSV.

A Lista de Adjacências é uma estrutura de dados composta por um vetor estático
de Nós (unidade básica de uma lista encadeada). No nosso caso, cada índice do
vetor representa um vértice diferente (v1, v2, ...), e cada lista associada a um
índice representa as adjacências do vértice correspondente.

\begin{figure}[h]
\centering
\begin{tabular}{c|l}
Vértice & Adjacentes \\ \hline
1 & 2, 3, 5 \\
2 & 1, 4 \\
3 & 1 \\
4 & 2, 5 \\
5 & 1, 4 \\
\end{tabular}
\caption{Representação de uma lista de adjacências. O vértice 1 é
conectado a 2, 3 e 5.}
\end{figure}

Foram implementadas funções que fornecem o valor máximo e mínimo de um vértice,
funções que determinam o tipo do grafo (simples ou multigrafo), além da quantidade
de laços e/ou arestas múltiplas que ocorrem nele e, como função principal, um
procedimento que determina quantos componentes conexos existem no grafo, bem como
sua distribuição e seus respectivos tamanhos.
\newpage
\section{Solução fornecida e sobre o projeto}

O projeto foi administrado em um \href{https://github.com/atomicpenguinz/TEG_1}{repositório}
do Github, e separado em arquivos diferentes para melhor organização do trabalho.
Um arquivo para as funções referentes à Lista de Adjacências, outro referente
a leitura e interpretação do arquivo CSV, outra para a interface de usuário e ainda
um arquivo header, que possui as definições e protótipos usados nos outros.

A seguir vai uma explicação simples sobre cada exigência:

\section{Grau mínimo e máximo}
\vspace{-0.3cm}
\rule{\textwidth}{0.4pt}
\vspace{0.3cm}


\newpage

\section{Descrição de atividades}

\newpage
\section{Análise dos Resultados}
\newpage
\section{Trabalhos Futuros}
\newpage

\addcontentsline{toc}{section}{Bibliografia}
\section*{Bibliografia}
\footnotesize{

\noindent AGUIRRE, L. A. Introdução à Identificação de Sistemas, Técnicas Lineares e Não lineares Aplicadas a Sistemas Reais. Belo Horizonte, Brasil, EDUFMG. 2004.\\

}
\newpage
\addcontentsline{toc}{section}{Anexo}
\section*{Anexo}
\end{document}




